% =============================================================================
\section{Prise en compte des aléas}
% =============================================================================

% Les quatre aléas.
\begin{frame}{Quatre types d'aléas à absorber}
  \small
  \renewcommand{\arraystretch}{1.3}
  \begin{tabular}{p{3.0cm}p{5.3cm}p{4.5cm}}
    \toprule
    \textbf{Aléa} & \textbf{Description} & \textbf{Impact opérationnel} \\
    \midrule
    Urgence & Patient imprévu, fenêtre temporelle imposée
      & À insérer dans une vacation compatible du jour \\
    Annulation & Patient programmé qui annule
      & Libère un créneau \textbf{et} des lits \\
    Indisponibilité lits & Réduction de capacité sur
      $[j_{\text{début}}, j_{\text{fin}}]$
      & Force la réaffectation des séjours chevauchants \\
    Retard bloc & Prolongation ou incident sur une vacation
      & Ampute la capacité utile disponible \\
    \bottomrule
  \end{tabular}

  \vspace{0.3cm}
  \begin{alertblock}{Deux réponses complémentaires}
    \small \textbf{Proactive} : générer des plannings alternatifs robustes.
    \textbf{Réactive} : adapter dynamiquement sous contrainte de stabilité.
  \end{alertblock}
  \source{Chaque aléa est modélisé par une classe dédiée et un simulateur
    reproductible (\texttt{generer\_scenario\_aleas}).}
\end{frame}

% Proactif : plannings alternatifs.
\begin{frame}{Proactif : plannings alternatifs et double proposition}
  \small
  \begin{columns}[T,onlytextwidth]
    \begin{column}{0.5\textwidth}
      \begin{block}{Quatre profils de planning}
        \begin{itemize}
          \item \textbf{Nominal} : optimal à 100\,\% de la capacité.
          \item \textbf{Robuste bufferisé} : réserve de capacité bloc
            ($\sim$15\,\%) et lits pour absorber les imprévus.
          \item \textbf{Sécurité lits} : lissage strict, sous-saturation pour
            parer aux fermetures.
          \item \textbf{Alternatif Date B} : haute qualité, diversifié par
            rapport au nominal.
        \end{itemize}
      \end{block}
    \end{column}
    \begin{column}{0.48\textwidth}
      \begin{exampleblock}{Double proposition chirurgicale}
        \footnotesize Pour chaque patient, la table
        \texttt{extraire\_options\_date\_a\_b} produit :
        \begin{center}
          \textbf{Date A} (optimale) \quad / \quad \textbf{Date B}
          (alternative)
        \end{center}
        Le praticien choisit en connaissance des tensions sur les lits et le
        bloc.
      \end{exampleblock}
      \begin{alertblock}{Objectif}
        \footnotesize Passer d'un planning \textbf{optimal} à un planning
        \textbf{robuste} : ne pas casser tout le programme parce qu'un imprévu
        survient.
      \end{alertblock}
    \end{column}
  \end{columns}
  \source{Aperçu : un buffer de capacité réduit la qualité nominale mais
    absorbe les aléas sans décaler les patients.}
\end{frame}

% Réactif : adaptation dynamique.
\begin{frame}{Réactif : adaptation sous contrainte de stabilité}
  \small
  \begin{columns}[T,onlytextwidth]
    \begin{column}{0.52\textwidth}
      \begin{block}{Trois garde-fous}
        \begin{itemize}
          \item \textbf{Gel du passé} : tout acte antérieur au jour courant est
            sanctuarisé.
          \item \textbf{Insertion prioritaire des urgences} dans leur fenêtre
            clinique.
          \item \textbf{Moindre perturbation} : on minimise les déplacements de
            patients futurs.
        \end{itemize}
      \end{block}
    \end{column}
    \begin{column}{0.46\textwidth}
      \begin{alertblock}{Score d'adaptation}
        \footnotesize
        \[
          \begin{aligned}
            \text{Score} = \text{Fitness} &- w_p \cdot \text{Déplacements} \\
                                          &- w_j \cdot \text{Changements de jour}
          \end{aligned}
        \]
        Le planning reste proche de celui annoncé aux patients.
      \end{alertblock}
      \begin{exampleblock}{Rapport d'arbitrage}
        \footnotesize Synthèse chiffrée des motifs de chaque déplacement et des
        urgences accueillies.
      \end{exampleblock}
    \end{column}
  \end{columns}
  \source{Rapport produit automatiquement par \texttt{adapter\_planning} pour
    garder la décision explicable.}
\end{frame}

% Resultat 1 : plan fige ou replanification (repris de Kaua_SMA_9_octobre.tex).
\begin{frame}{Résultats : plan figé ou replanification ?}
  \small
  \vspace{0.15cm}
  \centering
  \renewcommand{\arraystretch}{1.32}
  \setlength{\tabcolsep}{8pt}
  \begin{tabular}{llrr}
    \toprule
    Situation & Politique & Terminées / 279 & Dépassement (min) \\
    \midrule
    Sans fermeture & Figée ou réactive & 262 & 2\,954 \\
    \midrule
    Avec fermeture & Figée & \textbf{255} & 2\,911 \\
    Avec fermeture & Réactive & 251 & \textbf{2\,716} \\
    \bottomrule
  \end{tabular}
  \vspace{0.35cm}
  \begin{alertblock}{Compromis observé, pas un gain systématique}
    \small La replanification réduit le dépassement de \textbf{195 min}
    mais termine \textbf{4 interventions de moins} sur cette fenêtre.
  \end{alertblock}
  \source{Modèle prédictif, planificateur glouton, 42 lits synthétiques, 2 salles, 3--31 janvier 2022, graine 0. Dépassement = occupation + remise en état après 17h.}
\end{frame}

